\section{Background and Motivation}

The rapid integration of large language models (LLMs) into everyday life has created a new category of first-responder for personal health concerns. Surveys conducted in the United States document that a substantial and growing proportion of people turn to general-purpose LLMs specifically for mental health support, often as a first contact before, or instead of, seeking professional care \citep{rousmaniere2025mentalhealth}. In parallel, a growing literature evaluates the clinical quality and safety of LLM-generated mental health responses in high-resource settings \citep{hua2025mentalhealth, santos2026clinical, saha2025linguistic}. What this literature has not addressed systematically is whether the quality and character of those responses remain stable across the geographic and infrastructural diversity of the users who now rely on them.

This gap is not merely academic. Mental health services are unevenly distributed across the world: in many low- and middle-income countries, specialist services are scarce or absent, community awareness of crisis support options is limited, and dedicated crisis helplines either do not exist or are not widely documented \citep{osborn2024students}. For a user in such a setting, an LLM may function not as a supplementary resource but as the primary available channel for mental health guidance. If the quality, actionability, or cultural framing of that guidance varies systematically with where the user says they are, the models most likely to be a user's only option are also the most likely to provide inferior support.

A second, less studied problem concerns the measurement of such variation. Much of the existing literature on LLM quality and bias in healthcare relies on automated content coding --- keyword indicators, dictionary-based feature extraction, or embedding-distance metrics --- to convert free-text model output into quantifiable outcomes at scale \citep{dai2023llm, depaoli2024thematic}. These pipelines are rarely validated against independent human judgements before their outputs are used to support substantive claims. This dissertation demonstrates directly, through one of its own measurement failures, that automated content-coding measures can produce outputs that are internally consistent and inter-rule-reliable yet still measure the opposite of what they claim to measure, with the error detectable only through comparison with independently collected human labels. This methodological observation is offered as a contribution to the LLM-audit literature in its own right.

\section{Problem Statement}

This dissertation addresses two related problems. The primary problem is empirical: whether the geographic location a user discloses when posing a mental health query to an LLM systematically changes the actionability, localisation, and cultural framing of the guidance the model provides, and whether any such variation runs in the direction that would disadvantage users in lower-income and lower-infrastructure settings. The secondary problem is methodological: whether automated content-coding pipelines used to measure such variation can be trusted without independent validation, and what a validation process that is capable of catching failures --- including directional failures --- looks like in practice.

\section{Research Aim and Objectives}

The aim of this dissertation is to determine whether and how the geographic location implied by a user's stated city and country shapes the mental health guidance produced by contemporary large language models, using a methodology in which every automated outcome measure is independently validated before being used to support a hypothesis test.

The objectives are to construct a controlled, factorial dataset of LLM responses to a fixed battery of mental health scenarios, varying only the stated city and country of the user across twenty cities spanning a range of income levels and WHO regions; to develop automated text-based outcome measures aligned with three pre-registered hypothesis families; to validate each measure against independently collected human judgements, with pre-specified acceptability criteria, before its use in any confirmatory analysis; to test three pre-registered hypotheses, each with a single controlled statistical model, retaining city-level granularity throughout; and to assess robustness of any confirmed findings to the small number of geographic units and to the exclusion of any single model.

\section{Research Questions}

Three research questions are pre-registered. \textbf{RQ1} asks whether the actionability of LLM mental health responses -- the extent to which a response provides concrete, followable steps -- varies systematically with the income category of the user's stated location. \textbf{RQ2} asks whether the localisation of those responses -- the extent to which a response is grounded in the user's specific location rather than offering generic guidance -- varies systematically with the documented availability of a national crisis service in that location. \textbf{RQ3} asks whether the likelihood that an LLM recommends religious or spiritual support, as one form of cultural framing, varies systematically by WHO region. A fourth, exploratory finding -- the detection of visibly fabricated crisis-contact details -- emerges from auditing the raw response data and is reported separately from the confirmatory design.

\section{Scope and Contribution}

The study covers 1,120 responses from seven large language models across twenty cities, queried using a fixed battery of eight mental health scenarios that vary in acuity from moderate distress to presentations with an implied risk of self-harm. All queries are posed in English, holding language constant across locations; no real user data, patient records, or crisis disclosures are involved. The contribution is fourfold: an empirical answer to RQ1--RQ3 using validated measures; a documented case study in automated content-coding failure and its detection; a novel exploratory finding on the geographic distribution of fabricated crisis-contact information; and a fully reproducible pipeline released as electronic supplementary material.

\section{Dissertation Structure}

Chapter 2 reviews relevant literature on LLMs in medical and mental health contexts, geographic and cultural bias in language model outputs, and existing approaches to LLM auditing. Chapter 3 describes the dataset construction, the NLP pipeline, the validation methodology, and the statistical models. Chapter 4 presents the results, from exploratory data analysis through to confirmatory hypothesis tests and robustness checks. Chapter 5 discusses the findings in relation to the literature, their practical implications, and their limitations. Chapter 6 summarises the contributions and directions for future work.
